% Finite-scale producer for the neutral flat curvature representative.
\section{Finite-scale propagation of the neutral curvature representative}
\label{sec:ex-lattice-rg-construction}

This section records the finite-scale construction that supplies the scalar
recurrence used in Section~\ref{sec:ex-lattice-flat-rate}. It starts from a
terminal difference of two quadratic precisions, propagates a representative
of that difference modulo linear-curvature null directions, and controls its
cost at every earlier level. The one-step operation is the local Schur and
matching argument of Section~\ref{sec:ex-lattice-flat-step}; the present
section supplies its actual finite-history inputs and closes the resulting
recurrence. No limiting argument is used here.

\subsection{Data on a finite interval of levels}

For each $x$ in an index set $X$, fix one schedule, a measuring level
$k_x$, and finite indices $i_0(x)\leq N_x\leq n_x$. In formulas for one
fixed $x$, suppress its subscript and write
\[
 J_i=J(i),\qquad L_i=|\text{block at step }i|,\qquad
 t_i=\log a_i^{-1},\qquad \Omega_i=\Omega(J_i),\qquad
 \sigma_i=w(J_i),
\]
where $a_i$ is the positive length at level $i$. The schedule identities are
\[
 t_{i+1}-t_i=\log L_i,\qquad L_i\geq3.
\]
Let $\mathcal F$ be the finite set of face labels and put $m=|\mathcal F|$.
For a horizon metric $H$ at a torus side $c$, let $D_c$ be the discrete
exterior derivative from one-cochains to two-cochains and write its
quadratic precision as
\[
 \mathscr Q(H,c)=D_c^{\mathsf T}H_cD_c,
\]
where $H_c$ is the finite matrix obtained from the metric by periodization.
The packets below represent differences of these forms by periodizing a
cover curvature matrix between the two copies of $D_c$.
The scale allotment for adjacent levels and its one-step null correction are
\begin{equation}
 \mathcal A_i=m^2e^4\Omega_{i+1}^{2}\Omega_i,
 \qquad
 \nu_i(u)=\frac{2m^2\mathcal A_i^{2}}{uL_i}\quad(u>0).
 \label{eq:ex-lattice-rg-coin}
\end{equation}
These quantities are nonnegative. Indeed, the row bounds on each quadratic
precision imply $\Omega_i\geq0$. The expression $\mathcal A_i$ is the
allotment entering the two order-two response-kernel estimates. The factor
$m^2$ in $\nu_i$ counts the two finite face labels; the factor $L_i^{-1}$
and the factor $2/u$ are the centered null-packet correction.

Assume the following finite data.

\begin{enumerate}
\item For every member $x$ of an index set $X$, there is a horizon of
quadratic precisions $H_x(N,i)$, with row bounds at widths
$\alpha(J_i)$ and costs $\Omega_i$. These bounds include the diagonal,
crossing, and torus-floor estimates used in the one-step lemma. For every
$i_0(x)\leq i<N_x$, the actual Wilson tower precision at $(N_x,i)$ equals
the order-two precision of the flight from level $i+1$ to level $i$.
\item There is a positive width curve $w$, with
$w(k)\leq\alpha(k)$, such that for every $i<N$,
\begin{equation}
 \sigma_i\leq\min\{1,\alpha(J_i),\alpha(J_{i+1})\},\qquad
 2\sigma_i\leq L_i\sigma_{i+1},\qquad
 \nu_i(\sigma_{i+1})\leq
 \exp\bigl(-\theta(t_{i+1}-t_i)\bigr)
 \label{eq:ex-lattice-rg-funding}
\end{equation}
for one fixed $\theta>0$.
\item The terminal allowance
$T_{x,N}=2m\Omega_N$ has a uniform polynomial age bound
\begin{equation}
 T_{x,N}\leq A_0\,\mathfrak a_x(N)^d,
 \qquad
 \mathfrak a_x(i)=2+Z_x+|t_i-t_{k_x}|,
 \label{eq:ex-lattice-rg-topcap}
\end{equation}
with $Z_x\geq0$, integer $d\geq0$, and $A_0\geq0$ independent of
$x,N,n$.
\end{enumerate}

The first two items are the complete inputs represented by the plan and
funding conditions below. The first inequality in
\eqref{eq:ex-lattice-rg-funding} keeps all response kernels inside both
adjacent analytic widths. The second is the block-folding requirement. The
third pays the exact one-step null-packet cost.

\subsection{Production of the funded width and top allowance}
\label{subsec:ex-lattice-rg-funded-producer}

For a map price $X$, put
\[
 \omega_j=\min\{1,\alpha_X(j),\alpha_X(j+1)\},\qquad
 A_X(j,u)=m^2e^4\Omega_X(j)^2\Omega_X(u),
\]
\[
 b_X(j,u)=m A_X(j,u)\left(m\frac{2}{\omega_j}A_X(j,u)\right)
       =\frac{2m^2A_X(j,u)^2}{\omega_j}.
\]
\[
 E_X(j)=\left\lceil\frac{2}{\omega_j}+b_X(j,j)^2+b_X(j,j-1)^2\right\rceil.
\]
These are the source's entrance ticket and two adjacent row prices. The
schedule satisfies $J_{i+1}=J_i$ or $J_i+1$. Suppose the selected floors
give $E_X(J_{i+1})\leq L_i$ and the delay receipt gives
$\Omega_X(J_i)\leq\log\widehat\beta_i$. Then the world width is
$w(j)=\omega_j$ and $\theta=1/2$.

The width is positive, at most one, and no larger than $\alpha_X$. If the
stage rises, its bridge width is exactly $\omega_{J_i}$; if the stage stays
fixed, the bridge is $\min\{1,\alpha_X(J_i)\}\geq\omega_{J_i}$. The
entrance summand $2/\omega_{J_{i+1}}$ implies
$L_i\omega_{J_{i+1}}\geq2$, while $w(J_i)\leq1$, proving
$2w(J_i)\leq L_iw(J_{i+1})$.

Since $J_i$ is either $J_{i+1}$ or its predecessor,
$b_X(J_{i+1},J_i)^2\leq L_i$. The quote, by the definitions of allotment
and neutral coin, is exactly
\[
 \operatorname{quote}_i(w(J_{i+1}))
 =\frac{2m^2A_X(J_{i+1},J_i)^2}{w(J_{i+1})L_i}
 =\frac{b_X(J_{i+1},J_i)}{L_i}
 \leq L_i^{-1/2}
 =e^{-(t_{i+1}-t_i)/2}.
\]
This proves the paid condition and fixes its exponent from the actual
entrance-price formula.

For the top allowance let $s_G>0$ be the gauge slope constant and choose
$Z_x=3+\log\widehat\beta_{k_x}$. The selected spacing hypothesis includes $\widehat\beta_{k_x}>1$;
therefore $Z_x=3+\log\widehat\beta_{k_x}\geq0$. The schedule's
delayed-scale inequality and $\log r\leq r-1$ imply
$\log\widehat\beta_i\leq\log\widehat\beta_{k_x}
+s_G|t_i-t_{k_x}|$. Consequently
\[
 \Omega_X(J_i)\leq(1+s_G)(2+Z_x+|t_i-t_{k_x}|).
\]
As $T_{x,N}=2m\Omega_X(J_N)$, the top cap holds with degree one and
constant $2m(1+s_G)$. Thus the full funding conclusion follows from
the selected integer entrance floor and the post-length delay, rather than
being an independent large-scale assumption.

\subsection{Production of the horizon rows and exact tower identity}
\label{subsec:ex-lattice-rg-horizon-producer}

The selected public route supplies a horizon and an exact carrier identity,
but their local row construction is a separate producer obligation. The
source route retains finite-period rows, adds a Gram correction preserving
the curl action, takes square roots on a common smaller strip, and checks the
resulting block identities including harmonic modes. This paragraph records
that route and does not treat those operations as proved here. The exact
mathematical input consumed below is: for every $x,N,i$, a positive horizon
metric with the stated diagonal, crossing, floor and row bounds, and an
identity of its tower curvature with the order-two flight from $i+1$ to $i$
on every required finite torus. Producing these horizons and verifying the
Gram, root, and harmonic identities remains an open construction step in
this chapter. The native component recipe retains the same schedule when
binding these receipts to the selected run.
\subsection{The selected run's nonlinear one-step supplier}
\label{subsec:ex-lattice-volant-producer}

At the selected adjacent levels, let $L$ be the block side, $\beta$ the
lower scale, $B_0,B_1$ the lower and upper budget radii, $\mu$ the shape
strength, and $Q$ the menu target. The shipment tariff is
\begin{equation}
 N\,\operatorname{siteFactor}(B_0)
 \exp\!\left((Q.\mathrm{relative}+1)\operatorname{preCost}(Q.\mathrm{halo}+5)\right)a
 \leq \operatorname{ribbon}(Q,L,\beta),
 \label{eq:ex-lattice-volant-tariff}
\end{equation}
where $N$ is the occupancy count and $a$ is the available shipment amount.
The second-row allotment is
\begin{equation}
 \mathfrak s(\mu,L,\beta)
 =2\exp\!\left(-\frac{\mu}{16}\Delta_{G,L}(\beta)^2\right),
 \label{eq:ex-lattice-volant-straw}
\end{equation}
with $\Delta_{G,L}$ the fine-scale monitor.

The selected response produces ungrouped terms $A_i(U)$ and $D_i(U)$,
with support indices $s(i)$ and $p(i)$, respectively. Let $a_i,d_i\geq0$
be summable majorants for their weighted prices, and let $S_i,T_i$ be
their support sets. The precise paid inputs are
\begin{align}
 \operatorname{price}_r(A_i(U))&\leq a_i,&
 \sum_{i:z\in S_i}a_i e^{(b+1/2)\operatorname{coverNumber}(B_0,S_i)}&\leq1,
 \label{eq:ex-lattice-volant-first-ledger}\\
 \operatorname{price}_r(D_i(U))&\leq d_i,&
 \sum_{i:z\in T_i}d_i e^{(b+1)\operatorname{coverNumber}(B_0,T_i)}&\leq
       \mathfrak s(\mu,L,\beta).
 \label{eq:ex-lattice-volant-second-ledger}
\end{align}
Every $A_i$ has zero quiet coefficient and is supported in the assigned
support $S_i=s(i)$. Every $D_i$ is supported in the full box support
$T_i=p(i).1$ and in the prescribed threshold shell. Thus the cover weights
in the ledgers are the supports used by the resulting families themselves.
The selected response receipts
also give, for every box $p$ and configuration $U$,
\begin{align}
 \sum_{i:p(i)=p}[D_i(U)]_{\varnothing}&=\operatorname{profile}_p(U),
 \label{eq:ex-lattice-volant-quiet}\\
 \sum_{i:p(i)=p}[D_i(U)]_j&=\sum_l\operatorname{owl}(q_l)_j(p,U),
 \label{eq:ex-lattice-volant-field}
\end{align}
and the ungrouped expenditure identity
\begin{equation}
 \mathcal D(\mathcal E)(U)=e^{E(U)}\operatorname{pass}\!\left(\sum_i A_i(U)\right)
 *\operatorname{field}\!\left((\sum_{i:p(i)=p}D_i(U))_p,
                            \operatorname{flags}(U)\right).
 \label{eq:ex-lattice-volant-play}
\end{equation}
The field coefficient identity is for each box $p$ and singleton mode $j$;
higher coefficients are the full products in the field polynomial.

\begin{theorem}[Paid menu assembly and shipment conversion]
\label{thm:ex-lattice-volant-supplier}
Assume the selected history provides the ungrouped paid terms and exact
identities above, at lower exponent $b$ and upper exponent $e_1\geq100$.
Assume the menu offer has exponent $e_1-4$ and allowance $a$, and that the
birth estimate, foam, and radius hypotheses stated in the proof hold. Then
the grouped families
\[
 F_s(U)=\sum_{i:s(i)=s}A_i(U),\qquad
 K_p(U)=\sum_{i:p(i)=p}D_i(U)
\]
have the same rooted budgets as the ungrouped families, retain the support
and shell conditions, and satisfy the per-box quiet and singleton identities
and the push identity in \eqref{eq:ex-lattice-volant-quiet}--
\eqref{eq:ex-lattice-volant-play}. The selected door, tuning, and
slot-separation conditions are retained. The producers of the paid
terms are treated separately in Section~\ref{sec:ex-lattice-rg-multisource};
this theorem proves their tariff conversion and assembly.
\end{theorem}

\begin{proof}
For menu $i$, write $n_i^{\mathrm{birth}}$ for the cover number of its
birth region at radius $B_0$ and $n_i^{\mathrm{crest}}$ for the cover
number of its crest at radius $B_1$. The selected birth estimate is
\[
 (Q.\mathrm{relative}+1)n_i^{\mathrm{birth}}
 \leq (Q.\mathrm{relative}+1)\operatorname{preCost}(Q.\mathrm{halo}+5)
       +\frac54\varepsilon_Q n_i^{\mathrm{crest}}.
\]
Here $\varepsilon_Q>0$. Let $w_i$ be the menu wage, equal to its intrinsic
profit times its total menu fee, and set
$f_i=w_i\exp(2\varepsilon_Q n_i^{\mathrm{crest}})$. Foam at scale $L$,
$B_1\leq2B_0$, $Q.\mathrm{relative}\geq4$, and the profile hypotheses
give $4\varepsilon_Q<e_1$. Hence
\[
 \frac{13}{4}\varepsilon_Q<\frac{13}{16}e_1\leq e_1-4,
\]
where the last inequality follows from $e_1\geq100$. Adding the wage
weight to the birth estimate gives
\[
 e^{(Q.\mathrm{relative}+1)n_i^{\mathrm{birth}}} f_i
 \leq e^{(Q.\mathrm{relative}+1)\operatorname{preCost}(Q.\mathrm{halo}+5)}
       w_i e^{(e_1-4)n_i^{\mathrm{crest}}}.
\]
For every named marker $v$, the offer gives
$\sum_{i:v\in A_i}w_i e^{(e_1-4)n_i^{\mathrm{crest}}}\leq a$,
where $A_i$ is the set of available named origins for menu $i$. Choose, for each positive wage, a named origin in the
menu mode whose parent image lies in that menu's birth region. The birth
estimate and the displayed $13\varepsilon_Q/4$ margin transfer the offer
to a rooted budget on birth regions, with exponent
$Q.\mathrm{relative}+1$ and budget
$c=\exp((Q.\mathrm{relative}+1)\operatorname{preCost}(Q.\mathrm{halo}+5))a$.
Apply the rooted projection estimate: if the origin family has occupancy at
most $N$ at each parent site and every positive-wage mode has one such
origin in its birth support, then lowering the rooted exponent by one gives
\begin{equation}
 \sum_{i:z\in\mathrm{birth}_i}f_i
       e^{Q.\mathrm{relative}\,\operatorname{coverNumber}(B_0,\mathrm{birth}_i)}
 \leq N\operatorname{siteFactor}(B_0)c
 \leq\operatorname{ribbon}(Q,L,\beta).
 \label{eq:ex-lattice-volant-root-lift}
\end{equation}
The final inequality is exactly the purchased tariff
\eqref{eq:ex-lattice-volant-tariff}; the tariff is an input from the
schedule. Zero-wage terms contribute nothing.

Now group terms by their support and box indices. The first ledger already
uses exactly the supports $S_i=s(i)$, and the second uses exactly the box
supports $T_i=p(i).1$. By the triangle
inequality,
$\operatorname{price}_r(F_s(U))\leq\sum_{i:s(i)=s}a_i$ and
$\operatorname{price}_r(K_p(U))\leq\sum_{i:p(i)=p}d_i$. Summing these
bounds over fibers at a fixed site gives exactly the two budgets in
\eqref{eq:ex-lattice-volant-first-ledger}--
\eqref{eq:ex-lattice-volant-second-ledger}. Tonelli's theorem applies
because the majorants are nonnegative and summable. Absolute convergence
allows each finite marker coefficient to be grouped termwise, proving the
per-box quiet and singleton identities.

The pass is a finite coefficient operation and commutes with the absolutely
convergent grouping of $A_i$. The field operation is a finite sum of finite
products on the marker algebra and likewise commutes with grouping the
$D_i$ by box. Substituting both regroupings into the ungrouped expenditure
identity gives the stated push identity, including the exponential scalar
and the pass of $\sum_sF_s$. All support conditions pass to the unions
within each fiber. This proves the assembly without assuming the grouped
conclusion as an input.
\end{proof}

\subsection{The terminal representative}

\begin{lemma}[Terminal neutral representative]
\label{lem:ex-lattice-rg-terminal}
Let $H_x$ satisfy the row and floor bounds at level $N$, and suppose
$0<\sigma_N\leq\alpha(J_N)$. There is a summable, translation-covariant packet $B_N$ which is
neutral, has folded cost at most $T_{x,N}$ at width $\sigma_N$, and represents
\[
 \mathscr Q(H_x(n,N),c)-\mathscr Q(H_x(N,N),c)
\]
as a linear-curvature quadratic form on every torus side $c$ occurring at
level $N$.
\end{lemma}

\begin{proof}
Set $Y=H_x(n,N)$ and $Z=H_x(N,N)$, and define the packet by subtracting
their saddle kernels entrywise. The row bounds give the two floor
inequalities with the common parameter $\Omega_N$ and width $\sigma_N$.
For each metric row $R$, its saddle kernel is the plain translation-invariant
kernel evaluated on a constant output position. A strip bound of width
$\sigma_N$ and cost $\Omega_N$ makes this saddle packet's folded norm at
most $m\Omega_N$: sum first over the cover position and use the strip bound
for each of the $m$ output labels. The triangle inequality therefore bounds
the difference packet by $m(\Omega_N+\Omega_N)=T_{x,N}$. Translation
covariance follows from translation invariance of each plain kernel.

For neutrality, fix the root position and output face label and sum the
packet over the cover. Absolute convergence follows from the folded bound.
Each saddle sum is the zero-frequency row sum of its metric. The floor
condition pins both metric zero modes to the same identity matrix; hence
their row sums agree and the difference sum is zero. Finally, on each torus
side $c$, periodizing the packet gives the difference of the two metric
matrices. Each floor identifies its metric with the quadratic form obtained
from its full map, so subtracting the identities represents the precision
difference by the exterior derivative of the difference of those finite
curvature matrices. This proves the claim without inverting the curl
precision; the identity holds on the tested curvature form including its
harmonic zero mode.
\end{proof}

\subsection{One backward step and the actual kernel inputs}

\begin{lemma}[One-step predecessor]
\label{lem:ex-lattice-rg-one-step}
Fix $i<N$ and assume a packet $B_{i+1}$ represents the difference
$\mathscr Q(H_x(n,i+1),c)-\mathscr Q(H_x(N,i+1),c)$ on every required torus side. Suppose
the two order-two response kernels for the histories $n$ and $N$ have
common exponential width
\[
 \omega_i=\min\{1,\alpha(J_i),\alpha(J_{i+1})\}
\]
and cost $\mathcal A_i$, satisfy the city and matching identities with
the corresponding tower carriers, and the response kernel used for the
coarser history is a travelling kernel. Denote the finer-history response
by $T_i$ and the coarser-history response by $U_i$. If $B_{i+1}$ is neutral and
translation covariant and has folded cost $p\geq0$ at width
$u=\sigma_{i+1}$, then the composed packet
\[
 B_i=B_{i+1}\star(T_i,U_i)
\]
is neutral and translation covariant, represents the quadratic-precision
difference at level $i$ on every required torus side, and has folded cost
at width $\sigma_i$ at most $\nu_i(u)p$.
\end{lemma}

\begin{proof}
The two kernels are the responses obtained by composing the inverse
precision at the fine level $i+1$, the installed Wilson tower from $i+1$
to $i$, and the metric at the coarse level $i$, once for cutoff $n$ and
once for cutoff $N$. The deep
order-two tower estimate and the two adjacent precision row estimates give
their order-two atlas bounds. The priced atlas initially contains
$e^{4\omega_i}\Omega_{i+1}$ in its high-side entry; since
$0\leq\omega_i\leq1$, enlarging this to $e^4\Omega_{i+1}$ gives exactly
the allotment $\mathcal A_i$ in \eqref{eq:ex-lattice-rg-coin}.
The crossing identities for the actual towers give the required matches
with their periodized carrier maps. The response composition has the city
law because each factor has its own translation law and these laws add
under convolution. The widths are $\omega_i$ for both responses.

For the remaining travel condition, the Wilson tower's zero-alias column
is the toll-$L_i$ identity, namely $L_i^{-2}I$. Periodicity then gives a
position-independent column sum for the installed tower; this scalar
normalization is sufficient for the travel product. The fine cross-row bound
and the inverse-precision row bound make the two preceding convolutions
absolutely summable; absolute convergence permits changing their order.
The sum of a convolution of two travelling kernels is the product of their
position-independent row sums, so the full coarser response is travelling.
These are precisely the finite matching, city, and travel hypotheses in
Section~\ref{sec:ex-lattice-flat-step}; applying its proved Schur and
weighted-packet estimate gives the asserted identities and cost
$\nu_i(u)p$. The width, positivity, and folding assumptions in
\eqref{eq:ex-lattice-rg-funding} are exactly that lemma's remaining
hypotheses.
\end{proof}

\subsection{Backward induction and the scalar rate}

\begin{theorem}[Uniform finite-interval flat-sector bound]
\label{thm:ex-lattice-rg-construction}
Suppose the terminal and one-step hypotheses above hold for every
$x\in X$ and every finite interval $N\leq n$. There are nonnegative
constants $A_1$ and an integer $d_1$, and an exponent $c_1>0$, independent
of $x,N,n$, and packets $B_i$ for all $i\in[i_0(x),N_x]$, such that
each $B_i$ is neutral and translation covariant, represents
$\mathscr Q(H_x(n,i),c)-\mathscr Q(H_x(N,i),c)$ on all required torus sides, and has folded
cost
\[
 \|B_i\|_{\sigma_i,\mathrm{fold}}
 \leq A_1\,\mathfrak a_x(i)^{d_1}
       \left(\frac{a_N}{a_i}\right)^{c_1}.
 \]
\end{theorem}

\begin{proof}
Define the nonnegative allowance recursively by
\[
 C_N=T_{x,N},\qquad
 C_i=C_{i+1}\exp\bigl(-\theta(t_{i+1}-t_i)\bigr)
 \quad(i<N).
\]
At $i=N$, the terminal lemma constructs $B_N$ with cost at most $C_N$.
If $B_{i+1}$ has cost at most $C_{i+1}$, the one-step lemma and the
funding inequality give
\[
 \|B_i\|_{\sigma_i,\mathrm{fold}}
 \leq\nu_i(\sigma_{i+1})C_{i+1}
 \leq C_{i+1}\exp\bigl(-\theta(t_{i+1}-t_i)\bigr)=C_i.
\]
Neutrality, covariance, and the tested curvature identity are inherited
from the one-step lemma. Backward induction therefore constructs all
packets and proves their costs are bounded by $C_i$.

For each fixed $x$, the schedule times have gaps
$\Delta_i=t_{i+1}-t_i=\log L_i\geq\log3$. The scalar estimate of
Section~\ref{sec:ex-lattice-flat-rate}, applied on this finite interval
with terminal value $C_N$, zero interior source, source decay exponent
$c_0=\theta$, and the age
$R+Z_x+|t_i-t_{k_x}|$, converts the terminal polynomial bound
\eqref{eq:ex-lattice-rg-topcap} and the exact recurrence for $C_i$ into
\[
 C_i\leq A_1\bigl(2+Z_x+|t_i-t_{k_x}|\bigr)^{d_1}
       e^{-c_1(t_N-t_i)}.
\]
The scalar estimate gives $c_1=\theta/2$ and $d_1=d$. The fixed $R$ in
that estimate is absorbed into the age because
$2+Z_x+|t_i-t_{k_x}|\geq2$. The constants are independent of $x$ and
of the interval length. Since $e^{-(t_N-t_i)}=a_N/a_i$, this is the
claimed bound.
\end{proof}

\begin{remark}[Role in the selected construction]
\label{rem:ex-lattice-rg-public-route}
This finite-scale producer is the one used by the flat supply in the
selected construction. The public route first selects the Wilson run with
treasury and period funding and retains its original sheet, quote, and
post-length payments. Its native component choice includes the flat
entrance supplier. Realization applies that supplier to the very same
selected schedule and horizon. The flat entrance supplier obtains an
admissible spacing schedule and invokes the finite-history supplier whose
meshing step uses the world construction proved above. This chapter
discharges that row-level producer conditional on the stated funding,
response-kernel, tower-matching, and cost hypotheses. Their origin in the
complete public construction remains tracked in the accompanying audit;
the present theorem alone is not the final continuum or moment limit.
\end{remark}
